% SOURCE_PDF_PAGE: 155
会对这个结果感到不快。不过，作为计算机科学工作者，我们知道，不应该把这个数与精确的 0 比较，而应该在 \texttt{float64} 的精度范围内将它与 0 比较。下面这种写法可以检查 $\sin(\pi)$ 是否足够接近 0（达到 15 位数字的精度），从而可以认为二者没有区别：

\begin{gocode}[]
fmt.Println(math.Abs(math.Sin(math.Pi)-0) < math.Pow(10, -15))
\end{gocode}

\subsection{实现十进制数}

前面的理论内容可真不少。知道了这些，我们有多种方式可以实现这个 \texttt{Decimal} 结构体。我们选择把整数部分和小数部分分开。幸运的是，因为结构体的内容对包外不可见，用户并不依赖我们的具体实现，所以以后随时推翻这个决定，应该也不会破坏已经导出的 API。

这正是通常最好用自定义类型隐藏内部细节的另一个理由。实际开发时，我们本可以先用不精确的浮点数，以后再重构。不过我们还是别这么做了，因为我们已经知道浮点数会引入误差。

整数部分和小数部分是两个不同的数。但我们怎么知道这个小数在货币中表示什么呢？聪（satoshi）目前是比特币区块链所记录的最小货币单位，1 聪是一枚比特币的一亿分之一（0.00000001 BTC），这与欧元、法郎、格里夫纳或卢比通常采用的百分之一相去甚远。我们会把小数部分的精度保存为 10 的幂。这个精度值介于 0（我们总得能够表示 1.0）和一个不太大的数之间。因为我们不需要表示大于 $10^{30}$ 的数，所以不必

% SOURCE_PDF_PAGE: 156
保存大于 30 的 10 的指数。对于这么小的数，使用 \texttt{byte} 很稳妥。一个 \texttt{byte} 的最大值是 255，而我们肯定不需要那么大的 10 次幂。下面的代码清单展示了使用子单位和精度的实现。

\begin{gocode}[title={代码清单 6.6 \texttt{decimal.go}：\texttt{Decimal} 结构体的实现}]
// Decimal 可以表示具有固定精度的浮点数。
// 示例：
// 1.52 = 152 * 10^(-2) 将存储为 {152, 2} (1)
type Decimal struct {
    // subunits 是子单位的数量。
    // 将它乘以精度即可得到实际值
    subunits int64 // (2)
    // 一个单位中的“子单位”数量，以 10 的幂表示。
    precision byte // (3)
}
\end{gocode}

\begin{codeannotations}
\item 给出示例以帮助理解
\item 小数值
\item 小数值的精度，即 10 的幂的指数
\end{codeannotations}

现在，我们应该可以更新测试，加入真实的数量（\texttt{Decimal}）和货币（\texttt{Currency}）值了。但是，真的可以吗？

\subsubsection*{构造 \texttt{Decimal}}

结构体的字段没有导出，而我们确实想继续保持这种状态。我们需要为 \texttt{Decimal} 和 \texttt{Amount} 编写构造函数。

先从没有依赖项的 \texttt{Decimal} 开始。不过，它应该接收什么参数呢？如果要求提供三个 \texttt{int}，分别表示整数部分、小数部分和精度，那么以后就无法更改这种基于 \texttt{int} 的实现了。可以预期，为了从一开始就避免浮点数误差，调用方会把金额输入表示为字符串。

% SOURCE_PDF_PAGE: 157
在 Go 中，创建结构体的一种常见方式是使用 \texttt{New} 函数，正如第 4 章之前展示的那样。另一种方式是在大家都传递字符串时采用 \texttt{Parse} 命名模式，例如 \texttt{time.ParseDate} 或 \texttt{url.Parse}。我们的 \texttt{Decimal} 很适合采用这种模式。我们来编写一个函数，让它接收字符串参数，并返回该字符串所表示的 \texttt{Decimal}。

\subsubsection*{解析十进制数}

请在 \texttt{decimal.go} 文件中编写 \texttt{ParseDecimal} 函数，返回一个 \texttt{Decimal} 或一个错误。别忘了编写测试。你需要用 \texttt{strconv.ParseInt} 把字符串转换成整数，还需要用 \texttt{strings.Cut} 按分隔符切分字符串。在终端中，可以运行 \texttt{go doc strconv.ParseInt} 和 \texttt{go doc strings.Cut} 寻找一些灵感。请记住，我们要用 \texttt{int64} 表示数值，因为它是基本类型中容量最大的。下面简要列出了 \texttt{ParseDecimal} 需要完成的各个步骤：

\begin{gocode}[]
ParseDecimal(string) (Decimal, error)
// 1——找到 . 的位置，并以它为界切分。
// 2——把去掉 . 的字符串转换为整数。这一步可能失败
// 3——增加一些一致性检查
// 4——返回结果
\end{gocode}

把字符串 \texttt{"18.95"} 拆成 \texttt{"18"} 和 \texttt{"95"} 有好几种方法，\texttt{strings} 包提供了其中两种：\texttt{Cut} 和 \texttt{Split}。我们使用 \texttt{strings.Cut} 而不是 \texttt{strings.Split}，因为前者更简单，而且当输入字符串中没有分隔符时，用起来方便得多。

一方面，\texttt{Cut} 会在给定分隔符第一次出现之后，把字符串分成两部分，并返回一个布尔值，说明是否找到了分隔符。如果字符串不包含分隔符，该函数就返回完整字符串、一个空字符串和 \texttt{false}。

% SOURCE_PDF_PAGE: 158
另一方面，\texttt{Split} 会按分隔符把字符串拆成若干子串，并返回这些子串组成的切片。如果分隔符没有出现，\texttt{Split} 会返回只含一个字符串的切片。最后，如果分隔符或原字符串为空，它会返回一个空切片。

表 6.1 展示了 \texttt{strings.Cut} 和 \texttt{strings.Split} 处理三种不同字符串时的行为：分隔符没有出现、出现一次以及出现多次。

\medskip
\noindent\textbf{表 6.1 \texttt{strings.Cut} 与 \texttt{strings.Split} 的示例}

\begin{center}
\small
\begin{tabularx}{\textwidth}{|l|X|X|}
\hline
值 & \texttt{(strings.Cut("\{Value\}", "p")} & \texttt{(strings.Split("\{Value\}", "p")} \\
\hline
\texttt{banana} & \texttt{"banana", "", false} & \texttt{[]string\{"banana"\}} \\
\hline
\texttt{grape} & \texttt{"gra", "e", true} & \texttt{[]string\{"gra", "e"\}} \\
\hline
\texttt{apple} & \texttt{"a", "ple", true} & \texttt{[]string\{"a", "", "le"\}} \\
\hline
\end{tabularx}
\end{center}

从 \texttt{grape} 的 \texttt{strings.Split} 示例可以看到，结果切片的长度等于字母 ``p'' 的数量加 1。还值得注意的是，\texttt{strings.Split} 不会丢弃空字符串，这一点可以从 \texttt{apple} 示例看出来。由于 \texttt{ParseDecimal} 可以返回错误，我们先花点时间了解错误类型，以及如何检查它们。

\subsubsection*{错误类型}

解析这件事本身就几乎意味着可能会出现问题。如果用户传给我们的是字母，我们能怎

% SOURCE_PDF_PAGE: 159
么办？返回错误。\texttt{errors} 包提供了一个很有用的方法：

\begin{gocode}[]
errors.As(err error, target any) bool
\end{gocode}

它会报告 \texttt{err} 的具体值能否赋给 \texttt{target} 所指向的值。当你使用一个库，想把错误类型与领域错误作比较时，这会非常方便；例如，可以询问某个错误是否来自 \texttt{money} 包。当你是库的编写者时，为用户导出领域错误是一种体贴的做法，这样，如果错误来自你的库，他们就能在自己的代码中调整行为。

在这里，我们正是这群体贴的库作者，所以会在 \texttt{money} 包中导出一种领域错误类型，并实现标准 \texttt{errors} 包中的 \texttt{Error} 接口。我们来创建一种包专用的错误类型。

\begin{gocode}[title={代码清单 6.7 \texttt{errors.go}：\texttt{money} 包的自定义错误类型}]
package money

// Error 定义一种错误。
type Error string

// Error 实现 error 接口。
func (e Error) Error() string {
    return string(e)
}
\end{gocode}

这并没有花费多少力气，而且使用时的简洁性值得这点投入。最后，在开始写代码之前，我们来想想调用方能够理解哪些错误。如果待解析的字符串不是有效数字，就返回第一个错误。如果尝试处理过大的值，就触发第二个错误。设定上限是个好主意，因为它有助于确保不超过 \texttt{int64} 的最大值，尤其是在两个 \texttt{Decimal} 变量相乘时——这件事肯定会发生。

% SOURCE_PDF_PAGE: 160
下面的代码清单给出了这两个错误的声明。

\begin{gocode}[title={代码清单 6.8 \texttt{decimal.go}：导出错误}]
const (
    // 如果 decimal 格式错误，则返回 ErrInvalidDecimal。
    ErrInvalidDecimal = Error("unable to convert the decimal")

    // 如果数量过大，则返回 ErrTooLarge
    // 这会造成浮点精度错误。
    ErrTooLarge = Error("quantity over 10^12 is too large")
)
\end{gocode}

既然已经导出了能想到的错误，我们就来编写这个函数。函数代码见下面的代码清单。

\begin{gocode}[title={代码清单 6.9 \texttt{decimal.go}：\texttt{ParseDecimal} 函数}]
// maxDecimal 的值是一万亿，采用短级差——10^12。
const maxDecimal = 1e12

// ParseDecimal 把字符串转换为其 Decimal 表示。
// 它假定至多有一个小数分隔符，
// 并且分隔符是 '.'（英文句号字符）。
func ParseDecimal(value string) (Decimal, error) {
    intPart, fracPart, _ := strings.Cut(value, ".") // (1)

    subunits, err := strconv.ParseInt(
        intPart+fracPart, 10, 64) // (2)
    if err != nil {
        return Decimal{}, fmt.Errorf("%w: %s",
            ErrInvalidDecimal, err.Error()) // (3)
    }

    if subunits > maxDecimal {
        return Decimal{}, ErrTooLarge
    }

    precision := byte(len(fracPart)) // (4)
     return Decimal{subunits: subunits, precision: precision}, nil
}
\end{gocode}

\begin{codeannotations}
\item 切分整数部分和小数部分，并检查是否找到了后者
\item 解析该值
\item 如果无法解析，则返回错误
\item 精度就是小数分隔符之后给出的数字位数。
\end{codeannotations}

% SOURCE_PDF_PAGE: 161
\Needspace{14\baselineskip}
我们的 \texttt{precision} 变量是怎样工作的？来看表 6.2 中的几个例子。

\medskip
\noindent\textbf{表 6.2 精度示例}

\begin{center}
\begin{tabular}{|l|l|}
\hline
值（字符串） & 精度 \\
\hline
\texttt{5.23} & 2 \\
\hline
\texttt{2.15497} & 5 \\
\hline
\texttt{1} & 0 \\
\hline
\end{tabular}
\end{center}

可以看出，解析后数字的精度就是小数分隔符后的位数。如果用户给我们 1.1 美元，这有点奇怪，但仍然有效。请注意，把它转换为美元后会得到 \$1.10，多出一位数字，因为美元被划分到百分之一。

\subsubsection*{测试 \texttt{ParseDecimal}}

为了检查结果，我们需要访问 \texttt{Decimal} 结构体未导出的字段。要做到这一点，测试必须位于同一个包中，并了解具体实现。测试可以写成下面的代码清单这样。

% SOURCE_PDF_PAGE: 162
\begin{gocode}[title={代码清单 6.10 \texttt{decimal\_internal\_test.go}：测试 \texttt{ParseDecimal}}]
func TestParseDecimal(t *testing.T) {
    tt := map[string]struct {
        decimal string // (1)
        expected Decimal
        err      error
    }{
        "2 decimal digits": {
            decimal: "1.52",
            expected: Decimal{subunits: 152, precision: 2},
            err:       nil,
        },
        "no decimal digits": {...}, // (2)
        "suffix 0 as decimal digits": {...}, // (2)
        "prefix 0 as decimal digits": {...}, // (2)
        "multiple of 10": {...}, // (2)
        "invalid decimal part": {...},
        "not a number": { // (3)
            decimal: "NaN",
            err:    ErrInvalidDecimal,
        },
        "empty string": { // (3)
            decimal: "",
            err:    ErrInvalidDecimal,
        },
        "too large": {     // (3)
            decimal: "1234567890123",
            err:    ErrTooLarge,
        },
    }
    for name, tc := range tt {
        t.Run(name, func(t *testing.T) {
            got, err := ParseDecimal(tc.decimal)
            if !errors.Is(err, tc.err) {
                t.Errorf("expected error %v, got %v", tc.err, err)
            }
            if got != tc.expected {
                t.Errorf("expected %v, got %v", tc.expected, got)
            }
        })
    }
}
\end{gocode}

\begin{codeannotations}
\item 默认值：对代码没有用，但对读者有用
\item 一组常规测试用例
\item 一组错误测试用例
\end{codeannotations}

我们特别来看名为 \texttt{"suffix 0 as decimal digits"} 的测试。在这个测试中，我们解析数值 1.50，它会转换成一个拥有 150 个子单位、精度为 2 的 \texttt{Decimal}。这是正确的，但真的是我们能做到的最好结果吗？这里处理的是十进制数，所以保留这些额外的零没有意义，因为它们不提供任何信息。我们来使用 \texttt{Decimal} 类型的一个新

% SOURCE_PDF_PAGE: 163
方法 \texttt{simplify} 简化十进制数。这个方法将通过 \texttt{ParseDecimal} 的测试来检验。\texttt{simplify} 方法会不断删除最右侧的零，只要这么做不会影响 \texttt{Decimal} 的值；例如，32.0 应该简化为 32，但 320 必须保持为 320。

\begin{gocode}[title={代码清单 6.11 \texttt{decimal.go}：\texttt{simplify} 方法}]
func (d *Decimal) simplify() {
    // 对一个数使用 %10，会返回其十进制表示的最后一位。
    // 如果精度为正，这一位就属于
    // 小数分隔符右侧。
    for d.subunits%10 == 0 && d.precision > 0 {
        d.precision--
        d.subunits /= 10
    }
}
\end{gocode}

这是重要的第一步。记得定期运行测试并提交——提交信息要明确——尤其是在完成交付成果的某一部分时。处理完这段复杂逻辑后，编写 \texttt{Currency} 构造器就很容易了。

\subsection{\texttt{Currency} 值对象}

要理解这个输入数字是什么意思，我们需要一种货币。如前所述，每种货币都有固定精度，无法表示小于该精度的任何值：0.001 CAD 不是我们想表示的金额，因为现实生活中并不存在这种金额（它可能用于统计等非常特殊的场景，但不会用于交易）。我们没办法猜出每种货币的精度。可以通过服务获取这些信息，也可以读取数据库、读取文件等等；还可以简单地硬编码一份特殊货币清单，并把大多数货币采用的百分之一设为默认值。希望货币不会太频繁地开始采用新的子单位。毕竟，这只是一个业余项目，而且我们（暂时）并不打算支持稀奇古怪的历史货币。

% SOURCE_PDF_PAGE: 164
首先，需要把货币精度加入结构体，如下一份代码清单所示，它的取值范围是 0 到 3。\texttt{byte} 在这里同样是不错的选择。

\begin{gocode}[title={代码清单 6.12 \texttt{currency.go}：增加精度}]
// Currency 定义一种货币的代码
// 及其小数精度。(1)
type Currency struct {
    code      string
    precision byte // (2)
}
\end{gocode}

\begin{codeannotations}
\item 别忘了更新文档字符串。
\item 增加货币的精度
\end{codeannotations}

因为我们接收字符串并返回有效对象或错误，所以可以使用 \texttt{Parse} 前缀。我们来想想可能需要返回哪些错误。如果给出的货币代码无效，就应该能够返回错误。我们来新建一个 \texttt{Error} 类型的常量 \texttt{ErrInvalidCurrencyCode}，如代码清单 6.13 所示。可以看到，建议的错误名称以大写字母开头，表示它是导出的，这使包的调用方能够检查该错误。然后，可以创建一个名为 \texttt{ParseCurrency} 的函数，它接收字符串形式的货币代码，并返回一个 \texttt{Currency} 对象和一个错误。第一项验证是检查代码是否由三个字母组成；如果不是，可以直接返回这个新错误。否则，就对可能的货币代码进行 \texttt{switch}，返回具有相应精度的 \texttt{Currency} 对象。我们假定默认情况的精度是 2，因为大多数货币都有两位精度。

% SOURCE_PDF_PAGE: 165
\begin{gocode}[title={代码清单 6.13 \texttt{currency.go}：支持的货币所用的函数}]
// 当待解析货币不是标准的三字母代码时，
// 返回 ErrInvalidCurrencyCode。
const ErrInvalidCurrencyCode =
    Error("invalid currency code") // (1)

// ParseCurrency 返回名称所对应的货币，
// 并且可能返回 ErrInvalidCurrencyCode。
func ParseCurrency(code string) (Currency, error) {
    if len(code) != 3 { // (2)
        return Currency{}, ErrInvalidCurrencyCode
    }
    switch code {
    case "IRR":
        return Currency{code: code, precision: 0}, nil
    case "CNY", "VND": // (3)
        return Currency{code: code, precision: 1}, nil
    case "BHD", "IQD", "KWD", "LYD", "OMR", "TND": // (3)
        return Currency{code: code, precision: 3}, nil
    default:
        return Currency{code: code, precision: 2}, nil // (4)
    }
}
\end{gocode}

\begin{codeannotations}
\item 导出这个错误，以便调用方在这种情况下显示友好的消息
\item 按照 ISO-4217，检查货币是否由三个字母组成
\item Go 支持在同一行列出多个 \texttt{case}。
\item 按照 ISO-4217，其他所有流通货币都采用百分之一划分。
\end{codeannotations}

再次提醒，不要在生产环境中信任这个工具。在现实应用中，验证货币时应该对照一份可以更新、无须修改代码的清单。为了避免这个项目变得太大，装不进口袋，我们可以确保使用的是大写拉丁字母。

\begin{infobox}{支线任务 6.1}
确保货币代码由 A 到 Z 范围内的三个字母组成。如果想把事情弄复杂，可以使用 \texttt{regexp} 包；也可以逐一检查三个字节是否都位于 \texttt{'A'} 和 \texttt{'Z'} 之间。
\end{infobox}

% SOURCE_PDF_PAGE: 166
我们把一切都充分测试了吗？还没有。先试着自己为解析器编写测试，再看下面代码清单中我们的版本。

\begin{gocode}[title={代码清单 6.14 \texttt{currency\_internal\_test.go}：\texttt{TestParseCurrency} 函数},listing options={basicstyle=\ttfamily\footnotesize\color{PocketInk}}]
package money

import (
    "errors"
    "testing"
)


func TestParseCurrency_Success(t *testing.T) { // (1)
    tt := map[string]struct {
        in       string
        expected Currency
    }{
        "hundredth EUR": {
            in: "EUR",
            expected: Currency{code: "EUR", precision: 2}},
        "thousandth BHD": {...},
        "tenth VND": {...},
        "integer IRR": {...},
    }
    for name, tc := range tt {
        t.Run(name, func(t *testing.T) {
            got, err := ParseCurrency(tc.in)
            if err != nil {
                t.Errorf("expected no error, got %s", err.Error())
            }

             if got != tc.expected {
                 t.Errorf("expected %v, got %v", tc.expected, got)
             }
        })
    }
}


func TestParseCurrency_UnknownCurrency(t *testing.T) { // (2)
    _, err := ParseCurrency("INVALID")
    if !errors.Is(err, ErrInvalidCurrencyCode) {
        t.Errorf("expected error %s, got %v", ErrInvalidCurrencyCode, err)
    }
}
\end{gocode}

\begin{codeannotations}
\item 成功用例
\item 错误用例
\end{codeannotations}

这一次，我们没有使用验证函数，而是把成功用例与唯一的错误用例分开。这主要是个人偏好问题，因为一如既往，唯一的判断标准

% SOURCE_PDF_PAGE: 167
是下一位读者能否理解我们在测试什么，以及能否轻松添加或修改测试用例。调用 \texttt{t.Run} 还能确保所有测试用例都可以单独运行。如果测试通过，别忘了提交。现在，我们有了数值和货币，可以把二者组合起来，构成一笔 \texttt{Amount} 金额。

\subsection{NewAmount}

金额是以十进制数表示的某种货币数量。如前所述，如果一个十进制数的精度过高，它就可能与货币不兼容。例如，我们不应该允许创建 19.8875（精度为 4 的数）加拿大元，因为加拿大元的子单位是分（精度为 2）。构造一个无效对象毫无意义：\texttt{NewAmount} 函数的职责就是返回有效对象或错误。下面的代码清单展示了 \texttt{NewAmount} 函数的实现。

\begin{gocode}[title={代码清单 6.15 \texttt{amount.go}：\texttt{NewAmount} 函数}]
const (
    // 如果数字对货币而言
    // 精度过高，则返回 ErrTooPrecise。
    ErrTooPrecise = Error("quantity is too precise")
)

// NewAmount 返回一笔金额。
func NewAmount(quantity Decimal, currency Currency) (Amount, error) {
    if quantity.precision > currency.precision { // (1)
        // 为避免转换 0.00001 分，现在就退出。
        return Amount{}, ErrTooPrecise
    }
    quantity.precision = currency.precision // (2)

    return Amount{quantity: quantity, currency: currency}, nil
}
\end{gocode}

\begin{codeannotations}
\item 检查两个精度是否兼容
\item 确保使用货币的完整精度，而不是输入值的精度
\end{codeannotations}

% SOURCE_PDF_PAGE: 168
这个测试应该相当直截了当，因此这里不再给出我们的版本。当然，你仍然可以在本书的代码仓库中找到它。别忘了覆盖错误用例，并带覆盖率运行测试，确认没有遗漏任何内容：

\begin{shellcode}[]
> go test ./... -cover
\end{shellcode}

现在，在编写真正的转换之前，最后一步是更新 \texttt{Convert} 测试。这个测试只测试 \texttt{Convert}，并不测试 \texttt{ParseDecimal}、\texttt{ParseCurrency} 或 \texttt{NewAmount}（它们已经由各自的内部测试覆盖），那么该怎样处理潜在错误呢？检查 \texttt{ParseNumber} 及相关函数可能返回的各种值并不是 \texttt{TestConvert} 函数的职责，但它确实需要处理错误。

一种避免处理这些错误的办法，是编写一个函数，不做任何检查就构造所需结构，因为你知道测试用例是有效的。为了构造它们，你的函数需要位于 \texttt{money} 包中，并导出给 \texttt{money\_test} 包。可是这样一来，谁能阻止调用方使用这个测试工具函数，把无效值发送给你呢？这会让本节编写的真正构造器完全失去意义。我们不要选择这个危险的方案。另一种办法是调用函数并忽略错误：

\begin{gocode}[]
number, _ := money.ParseDecimal("23.52")
\end{gocode}

这种写法很实用，也没有导出任何危险的东西，但它带来了一个问题：如果不小心给测试提供了无效值，就完全检测不到。我们会在自己毫不知情的情况下，对函数随错误一并返回的零值进行测试，从而产生不可靠的测试。我们同样不要选择这个方案。

% SOURCE_PDF_PAGE: 169
由于需要依次调用一连串不同的构造器，我们改为编写一个测试辅助函数，并把这一点告诉 Go（见代码清单 6.16）。用于单元测试的 \texttt{testing.T} 对象拥有一个 \texttt{Helper} 函数，其文档（运行 \texttt{go doc testing.T.Helper} 查看）这样写道：“\texttt{Helper} 会把调用函数标记为测试辅助函数。打印文件和行信息时，将跳过该函数。”这意味着你看到的是哪项测试出错，而不是这个辅助函数的行号。

% SOURCE_PDF_PAGE: 170
\begin{gocode}[title={代码清单 6.16 \texttt{convert\_test.go}：解析 \texttt{Currency} 和 \texttt{Amount} 的辅助函数}]
package money_test

import (
    "testing"

    "learngo-pockets/moneyconverter/money"
)

func mustParseCurrency(t *testing.T, code string) money.Currency {
    t.Helper() // (1)

    currency, err := money.ParseCurrency(code)
    if err != nil {
        t.Fatalf("cannot parse currency %s code", code) // (2)
    }

    return currency
}

func mustParseAmount(t *testing.T, value string, code string) money.Amount {
    t.Helper() // (3)

    n, err := money.ParseDecimal(value)
    if err != nil {
        t.Fatalf("invalid number: %s", value)
    }

    currency, err := money.ParseCurrency(code)
    if err != nil {
        t.Fatalf("invalid currency code: %s", code)
    }

    amount, err := money.NewAmount(n, currency)
    if err != nil {
        t.Fatalf("cannot create amount with value %v and currency code %s",
            n, code)
    }

    return amount
}
\end{gocode}

\begin{codeannotations}
\item 在调用栈中跳过这个函数
\item 不只是让测试失败，而是就在这里彻底停止
\item 在调用栈中跳过这个函数
\end{codeannotations}

可以看到，我们没有使用 \texttt{t.Fail}，而是用了 \texttt{t.Fatal}，后者会立即停止测试运行。现在可以给 \texttt{Convert} 函数的测试提供实际值了，如下一份代码清单所示。不过，返回值仍然是零值。

% SOURCE_PDF_PAGE: 171
\begin{gocode}[title={代码清单 6.17 \texttt{convert\_test.go}：在测试用例中调用 \texttt{mustParse}}]
"34.98 USD to EUR": {
    amount:   mustParseAmount(t, "34.98", "USD"), // (1)
    to:       mustParseCurrency(t, "EUR"), // (2)
    validate: func(t *testing.T, got money.Amount, err error) {
        if err != nil {
            t.Errorf("expected no error, got %s", err.Error())
        }
        expected := money.Amount{}
        if !reflect.DeepEqual(got, expected) {
            t.Errorf("expected %q, got %q", expected, got)
        }
    },
},
\end{gocode}

\begin{codeannotations}
\item 输入金额
\item 输入货币
\end{codeannotations}

能编译吗？能运行吗？测试能通过吗？干得好！现在正适合提交你的工作。

\section{应用转换逻辑}

我们对这个包的 API 感到满意（或者至少觉得还行），而且手中的对象都保证有效并且受支持。现在，真正让它转换金额吧。在第一个版本中，在能够运行工具之前，我们先硬编码一个汇率。验证基础逻辑后，就可以调用掌握真实数据的远程服务器。我们决定使用欧洲中央银行作为这个远程服务器；它提供免费服务，不需要身份验证，而且几年后很可能仍然存在。

\subsection{应用汇率}

我们创建的所有实体中，哪一个应该负责应用汇率？这个逻辑可以属于 \texttt{Amount} 结构体，因为它知道怎样用新值创建新的 \texttt{Amount}。当然，金额应该不可变，而且必须确保这项操作不会修改

% SOURCE_PDF_PAGE: 172
输入金额。但是，从概念上说，这个选择合理吗？你会期待给定货币的一笔钱告诉你它值多少另一种货币吗？你会期待手里的 10 英镑纸币告诉你：“嘿，我今天大约值 10 美元”吗？大概不会。你会去兑换处，递出纸币，并希望拿回另一张纸币和几枚硬币。如果概念上就不合理，未来的维护者也会更难理解。因此，我们把兑换处写成一个由 \texttt{Convert} 调用的函数。

\subsubsection*{实现 \texttt{applyExchangeRate}}

如前所述，这部分逻辑不打算使用 \texttt{float64}。它最为敏感，我们希望保证计算完全精确，不丢失所处理数值的任何精度。一边是 \texttt{Amount} 的 \texttt{quantity} 字段，类型为 \texttt{Decimal}；另一边是通过远程调用获取的汇率。还必须提供目标货币，因为输出金额的精度就存放在那里。

应该怎样表示这个汇率呢？欧洲中央银行发布的汇率最多有七位数，这意味着可以安全地把它存进 \texttt{float64} 变量。对于使用七位数字的货币，\texttt{float32} 可能不够（谁知道为什么不会再增加第八位）。我们已经创建了 \texttt{Decimal}，把它作为高精度浮点数的专用类型，所以最好使用它。由于这个变量不会表示“普通”十进制数，不妨采用一个专门类型，以最恰当地描述它的用途。甚至可以为它创建一个新文件；不过到了这里，就连这些作者当中最坚定的小文件拥护者也会承认，这个类型也可以

% SOURCE_PDF_PAGE: 173
放在 \texttt{convert.go} 文件里，紧跟在导出方法之后（见下面的代码清单）。

\begin{infobox}{注意}
在一个文件中，你总会希望先放导出方法，因为从代码的入口点开始阅读更容易。如果同一文件中有多个导出函数，可以先放一个导出函数，再把它调用的未导出函数紧随其后。
\end{infobox}

\begin{gocode}[title={代码清单 6.18 \texttt{convert.go}：\texttt{ExchangeRate} 类型}]
// ExchangeRate 表示从一种货币转换为另一种货币的汇率。
type ExchangeRate Decimal
\end{gocode}

这就为函数带来了明确的签名：\texttt{applyExchangeRate} 负责用汇率乘以输入数量，并返回一个与目标 \texttt{Currency} 兼容的 \texttt{Amount}。首先需要一个函数，将 \texttt{Decimal} 与 \texttt{ExchangeRate} 相乘，然后还要调整乘积的精度，使之与 \texttt{Currency} 的精度相符。

为了让十进制数与汇率相乘，我们把汇率转换成 \texttt{Decimal} 并做一些数学运算：乘法结果的值是两个值的乘积，而返回的十进制数精度是两个精度之和。过程如下：

\begin{gocode}[]
20.00*4.0 = {subunits: 2000, precision: 2}*{subunits: 40, precision: 1}
          = {subunits: 2000*40, precision: 2+1}
          = {subunits: 80_000, precision: 3}
\end{gocode}

第一点需要注意：我们用了远多于必要位数的数字，最终得到 80。尽管 80 等于 80.000，但我们其实并不需要这种精度。执行乘法时，可以再次使用 \texttt{simplify} 方法。

% SOURCE_PDF_PAGE: 174
第二点需要注意：我们目前得到的精度还没有考虑任何货币信息。到现在为止，所做的只是十进制数乘法。因此，在 \texttt{applyExchangeRate} 中，需要调整 \texttt{multiply} 的结果，使其具有目标货币的精度。为此，必须根据目标货币精度与汇率乘积精度之差进行乘法（或除法）。当然，可以在这里直接调用 \texttt{math.Pow(10., precisionDelta)}，但这样成本很高，需要在浮点数和整数之间进行大量转换。我们改为把这项任务委托给名为 \texttt{pow10} 的函数。在这个函数中，先硬编码几个常用的 10 的幂作为快速路径；只有指数超出预期范围时，才默认采用代价高昂的 \texttt{math.Pow} 调用，如代码清单 6.19 所示。总的来说，这个 \texttt{pow10} 函数既可以用包含所有情况的映射实现，也可以用 \texttt{switch} 语句实现。我们选择后者，但两种做法都可以。

\begin{gocode}[title={代码清单 6.19 \texttt{decimal.go}：\texttt{pow10()}}]
// pow10 是把 10 提升到给定幂次的一种快速实现。
// 它针对较小的幂做了优化，而对异常高的幂会比较慢。
func pow10(power byte) int64 {
    switch power {
    case 0:
        return 1
    case 1:
        return 10
    case 2:
        return 100
    case 3:
        return 1000
    default:
        return int64(math.Pow(10, float64(power))) // (1)
    }
}
\end{gocode}

\begin{codeannotations}
\item 调用 \texttt{math.Pow} 比直接返回一个值慢得多。
\end{codeannotations}

我们先为这第一部分编写代码，然后就能实现那个神秘的 \texttt{multiply} 函数，如代码清单 6.20 所示。这里的 \texttt{switch} 用来根据目标货币的精度调整结果。

% SOURCE_PDF_PAGE: 175
\begin{gocode}[title={代码清单 6.20 \texttt{convert.go}：\texttt{applyExchangeRate}}]
// applyExchangeRate 返回一个新 Amount，表示输入值
// 乘以汇率后的结果。
// 返回值的精度就是目标 Currency 的精度。
// 此函数不保证输出金额受支持。
func applyExchangeRate(a Amount, target Currency, rate ExchangeRate) Amount {
    converted, err := multiply(a.quantity, rate) // (1)
    if err != nil {
        return Amount{}
    }

    switch { // (2)
    case converted.precision > target.precision: // (3)
        converted.subunits =
            converted.subunits / pow10(converted.precision-target.precision)
    case converted.precision < target.precision: // (4)
        converted.subunits =
            converted.subunits * pow10(target.precision-converted.precision)
    }

    converted.precision = target.precision

    return Amount{
        currency: target,
        quantity: converted,
    }
}
\end{gocode}

\begin{codeannotations}
\item 把汇率应用到输入数量
\item 检查两个精度
\item 需要截掉一些数字。
\item 需要追加一些零。
\end{codeannotations}

返回的 \texttt{Amount} 不是用负责验证的函数构造的。相反，我们选择返回一个金额，让 \texttt{Convert} 函数在把它交给外部调用方之前负责验证。请注意，我们在文档中明确说明了这一点：如果未来的维护者（也包括你）出于某种原因想导出这个函数，就必须修改它，使其在必要时返回错误。

最后，本章的核心就在乘法函数中，来实现它吧！请记住，我们不想让浮点数彼此相乘，因为这可能导致浮点错误。这意味着必须把 \texttt{ExchangeRate} 转换成 \texttt{Decimal}，如下一份代码清单所示。其余部分相当直截了当。

% SOURCE_PDF_PAGE: 176
\begin{gocode}[title={代码清单 6.21 \texttt{convert.go}：\texttt{multiply}}]
// 将 Decimal 与 ExchangeRate 相乘并返回乘积
func multiply(d Decimal, r ExchangeRate) Decimal {
    dec := Decimal{
        subunits: d.subunits * rate.subunits,
        precision: d.precision + rate.precision,
    }
    // 稍微整理一下表示形式。删除末尾的零。
    dec.simplify()

    return dec
}
\end{gocode}

现在有了把金额转换成新货币的函数，应该在哪里调用它？回答这个问题之前，先测试一下刚写的内容。

\Needspace{0.78\textheight}
\subsubsection*{测试 \texttt{applyExchangeRate}}

这是逻辑的核心。需要进行大量测试，确保一切正常工作；即便以后决定更改某种实现，也能继续正常工作。

编写测试之前，可以先想想所有可能的测试用例。表 6.3 给出了几个例子。

% SOURCE_PDF_PAGE: 177
\medskip
\noindent\textbf{表 6.3 可能的测试用例}

\begin{center}
\footnotesize
\begin{tabularx}{\textwidth}{|l|l|l|X|}
\hline
金额 & 汇率 & 目标货币精度 & 我们要检查的内容 \\
\hline
\texttt{1.52} & \texttt{1} & \texttt{2} & 输出必须与输入完全相同。 \\
\hline
\texttt{2.50} & \texttt{4} & \texttt{2} & 小数部分变为 0，但精度应该是目标货币的精度。 \\
\hline
\texttt{4} & \texttt{2.5} & \texttt{0} & 小数部分变为 0，精度也是 0。 \\
\hline
\texttt{3.14} & \texttt{2.52678} & \texttt{2} & 真实汇率 \\
\hline
\texttt{1.1} & \texttt{10} & \texttt{1} & 保持精度 1 \\
\hline
\texttt{1\_000\_000\_000.01} & \texttt{2} & \texttt{2} & 在大数中保持精度 \\
\hline
\texttt{265\_413.87} & \texttt{5.05935e-5} & \texttt{2} & 非常小的汇率 \\
\hline
\texttt{265\_413} & \texttt{1} & \texttt{3} & 输入没有精度时，在输出中增加额外精度 \\
\hline
\texttt{2} & \texttt{1.337} & \texttt{5} & 提高输出精度 \\
\hline
\texttt{2} & $1.33 * 10^{16}$ & \texttt{5} & 汇率过高。 \\
\hline
\end{tabularx}
\end{center}

% SOURCE_PDF_PAGE: 178
不同测试用例的数量，以及我们不断想起新边界情况的速度，都表明应该采用基于表格或映射的测试。当然，由于这个函数没有导出，测试必须是内部测试，也就需要新建一个文件。测试实现可以写成下面的代码清单这样。

% SOURCE_PDF_PAGE: 179
\begin{gocode}[title={代码清单 6.22 \texttt{convert\_internal\_test.go}：测试 \texttt{applyExchangeRate}}]
package money

import (
    "reflect"
    "testing"
)

func TestApplyExchangeRate(t *testing.T) {
    tt := map[string]struct {
        in              Amount
        rate            ExchangeRate
        targetCurrency Currency
        expected        Amount
    }{
        "Amount(1.52) * rate(1)": { // (1)
             in: Amount{
                 quantity: Decimal{
                     subunits: 152,
                     precision: 2,
                 },
                 currency: Currency{code: "TST", precision: 2},
             },
             rate:           ExchangeRate{subunits: 1, precision: 0},
             targetCurrency: Currency{code: "TRG", precision: 4},
             expected: Amount{
                 quantity: Decimal{
                     subunits: 15200,
                     precision: 4,
                 },
                 currency: Currency{code: "TRG", precision: 4},
             },
        },

        // 添加测试用例
    }

    for name, tc := range tt {
        t.Run(name, func(t *testing.T) {
            got := applyExchangeRate(tc.in, tc.targetCurrency, tc.rate)
            if !reflect.DeepEqual(got.number, tc.expected) {
                t.Errorf("expected %v, got %v", tc.expected, got)
            }
        })
    }
}
\end{gocode}

\begin{codeannotations}
\item 测试用例用参数命名。
\end{codeannotations}

最后，\texttt{Convert} 现在可以向调用方返回有用的内容了。目前还没有汇率，所以先把汇率硬编码为 2，并把获取汇率的工作留到本章后面。

% SOURCE_PDF_PAGE: 180
\begin{gocode}[title={代码清单 6.23 \texttt{convert.go}：\texttt{Convert} 的第一个实现}]
// Convert 应用汇率，把一笔金额转换为目标货币。
func Convert(amount got. Number, to Currency) (Amount, error) {
    // 应用获取到的汇率，转换为目标货币。
    convertedValue := applyExchangeRate(amount, to, 2) // (1)
    return convertedValue, nil // (2)
}
\end{gocode}

\begin{codeannotations}
\item 调用应用汇率的函数
\item 返回新值
\end{codeannotations}

恭喜，你把这个函数的测试弄坏了！现在我们会返回一些内容，所以你可以调用 \texttt{mustParseAmount} 来定义预期输出，从而修复测试。既然我们已经信任转换的核心，就可以确保返回给调用方的内容还能再次被自己的库使用。

\subsection{验证结果}

由于支持的值有这么多限制，检查能否得到预期输出似乎很明智，如代码清单 6.24 所示。限制如下：

\begin{itemize}
\item 数值不能高于 $10^{12}$，否则会由于浮点数而丢失精度。
\item 数值的精度必须与货币的精度兼容。
\end{itemize}

与转换逻辑不同，这项工作可以由 \texttt{Amount} 结构体负责。一笔金额可能有效，也可能无效，而 \texttt{Amount} 应该知道自己是否有效。毕竟，英镑纸币和美元钞票上布满了必需的华丽装饰，用来证明其真实性。

% SOURCE_PDF_PAGE: 181
\begin{gocode}[title={代码清单 6.24 \texttt{amount.go}：\texttt{validate} 方法的实现}]
// 当且仅当 Amount 使用起来不安全时，validate 才返回错误。
func (a Amount) validate() error {
    switch {
    case a.quantity.subunits > maxDecimal: // (1)
        return ErrTooLarge // (1)
    case a.quantity.precision > a.currency.precision:
        return ErrTooPrecise // (1)
    }

    return nil
}
\end{gocode}

\begin{codeannotations}
\item 常量和错误是在解析数字时定义的。
\end{codeannotations}

代码清单 6.25 展示了 \texttt{Convert} 函数最终的样子。它非常短小，所以内部需要测试的内容不多。

\begin{gocode}[title={代码清单 6.25 \texttt{convert.go}：\texttt{Convert} 的实现}]
// Convert 应用汇率，把一笔金额转换为目标货币。
func Convert(amount a.quantity. Subunits, to Currency) (Amount, error) {
    // 应用获取到的汇率，转换为目标货币。
    convertedValue := applyExchangeRate(amount, to, 2)

    // 验证转换后的金额是否处于可处理的有界范围内。
    if err := convertedValue.validate(); err != nil {
        return Amount{}, err
    }

    return convertedValue, nil
}
\end{gocode}

检查一下测试：最终完成的库有没有令人信服的覆盖率？可以查看测试覆盖率。

现在，库的整个结构和逻辑都已经完成并经过测试，让我们把它接起来吧，因为代码不能运行就不好玩了！之后，我们会获取一些现实中的汇率并完成这个工具。拥有一个不断增加功能的可执行程序，就可以展示产品的早期版本，再逐步加以改进。

% SOURCE_PDF_PAGE: 182
\section{编写 CLI}

本节将编写 \texttt{main} 函数：不要忘记，我们不是在编写库，而是在编写 CLI。这意味着要解析并验证输入参数，然后把它们传给 \texttt{Convert} 函数。不过，在实现所有这些安全网之前，我们想先让应用程序运行起来。

\subsection{标志和参数}

退一步想想。程序应该做什么？来看本章开头写下的需求用法：

\begin{shellcode}[]
> change -from EUR -to USD 413.98
\end{shellcode}

请注意，这里有两个标志（输入货币和输出货币）以及一个参数（想兑换的金额）。

\subsubsection*{货币标志}

要从命令行读取标志，正如第 2 章所见，Go 有一个含义明确的 \texttt{flag} 包。在 \texttt{main.go} 文件中导入它后，可以先写最前面的几行代码，确保正确读取命令行。

如第 2 章所述，\texttt{flag} 包提供了一些很有用的方法，可以获取标志参数中的值。在代码清单 6.26 中，我们会使用 \texttt{flag.String} 方法，获取源货币和目标货币两个值：\texttt{from} 和 \texttt{to}。

\texttt{flag.String} 方法接收标志名称、默认值（可以为空）和简短说明作为参数。它以 \texttt{*string} 类型的变量返回 \texttt{-from} 标志的内容。正如第 2 章已经提到的，在定义所有

% SOURCE_PDF_PAGE: 183
标志之后、访问值之前，必须调用 \texttt{Parse} 方法。这里，我们把 \texttt{-from} 标志的默认值留空，却把 \texttt{-to} 标志的默认值设成字符串 \texttt{EUR}。如果用户没有在命令行中提供 \texttt{-from} 标志，\texttt{from} 变量的值就是空字符串。同样，如果没有 \texttt{-to} 标志，它的值就是 \texttt{EUR}。

\begin{gocode}[title={代码清单 6.26 \texttt{main.go}：解析标志}]
package main

import (
    "flag"
    "fmt"
)

func main() {
    from := flag.String("from", "", "source currency, required")
    to := flag.String("to", "EUR", "target currency") // (1)

    flag.Parse()

    fmt.Println(*from, *to) // (2)
}
\end{gocode}

\begin{codeannotations}
\item 可以使用本地货币作为默认值。
\item 打印出来，检查效果
\end{codeannotations}

现在可以运行它了。经历这么多库开发之后，你还记得怎样从终端运行程序吗？

\begin{shellcode}[]
> go run . -from EUR -to CHF 10.50
\end{shellcode}

\texttt{main} 函数的第一个实现没有调用 \texttt{Convert} 函数，但确实会打印源货币和目标货币。屏幕上应该会显示它们。

\subsubsection*{金额参数}

实现 CLI 的下一步是获取需要转换的值。如果命令行中没有这个值，我们就带错误退出。

% SOURCE_PDF_PAGE: 184
\subsubsection*{从命令行获取参数}

运行可执行程序时，大多数时候都需要指定输入、行为、输出等等。这些参数可以通过命令行显式提供，也可以通过预设环境变量或位于已知位置的配置文件隐式提供。对于显式设置，可以用两种方式把用户定义的值传给程序：参数和标志。

参数是一个从 0 开始、依次排列的参数序列。参数没有名称，而且有顺序。交换它们的位置，程序的行为可能完全改变。大多数时候，参数是必需的，也没有默认值。\texttt{go build} 命令唯一的参数，是包含 \texttt{main} 函数的目录、文件或包。

另一方面，标志参数不受位置顺序约束。它们可以在命令行中以任意顺序出现，而不会改变程序的行为。它们可以拥有默认值（命令行中没有该标志时使用），就像我们的 \texttt{-to} 一样。你可能用过一个控制行为的标志例子：\texttt{go build} 命令的 \texttt{-o \{binaryPath\}} 选项。

在 Go 中，可以通过 \texttt{os.Args} 或 \texttt{flag} 包获取命令行参数。下面的例子展示了二者的差异：

\begin{shellcode}[]
> ./convert -from EUR -to JPY 15.23
\end{shellcode}

在这一行中，\texttt{os.Args} 返回一个含六个字符串的列表，每一项代表命令中的一个词：\texttt{\{"./convert", "-from", "EUR", "-to", "JPY", "15.23"\}}。而 \texttt{flag.Args} 则

% SOURCE_PDF_PAGE: 185
只返回命令行中不属于标志的参数：\texttt{\{"15.23"\}}。

该使用 \texttt{flag.Args} 还是 \texttt{os.Args}，取决于我们想访问什么信息。这里，我们只想访问命令行参数，并不关心设置了哪些标志。使用第一个不属于标志的参数非常简单：可以使用 \texttt{flag.Arg(0)}。

\begin{gocode}[title={代码清单 6.27 \texttt{main.go}：获取程序的第一个参数}]
func main() {
    from := flag.String("from", "", "source currency, required")
    to := flag.String("to", "EUR", "target currency")

    flag.Parse()

    value := flag.Arg(0) // (1)
    if value == "" {
        _, _ = fmt.Fprintln( // (2)
            os.Stderr, "missing amount to convert")
        flag.Usage() // (3)
        os.Exit(1)
    }

    fmt.Println(*from, *to, value) // (4)
}
\end{gocode}

\begin{codeannotations}
\item 如果没有参数，则返回空字符串
\item 忽略 \texttt{Fprintln} 的输出
\item 打印错误和工具的用法
\item 打印出来，检查是否读到了所有内容
\end{codeannotations}

输入已经拿到了。它们是字符串，而我们并不确定值是否有效。幸运的是，我们已经有了恰好适合处理它们的函数。

\subsection{解析为业务类型}

\texttt{Convert} 函数接收的参数值已经具有适合它使用的类型，而包也导出了构造这些值的方法。这种策略让调用方的逻辑尽可能灵活，因为 \texttt{main} 可以让类型经过自己新增的其他任何逻辑，也可以使用自己的不同类型，然后

% SOURCE_PDF_PAGE: 186
在最后一刻调用 \texttt{Parse}；还可以携带字符串，在每次需要时调用 \texttt{Parse}。

本章除了转换以外不会再做太多事情（欢迎你以后自行扩展）。我们只需要把它们全部解析出来。

% SOURCE_PDF_PAGE: 187
\begin{gocode}[title={代码清单 6.28 \texttt{main.go}：解析 \texttt{Currency} 和 \texttt{Amount}}]
package main

import (
    "flag"
    "fmt"
    "os"

    "learngo-pockets/moneyconverter/money"
)

func main() {
    from := flag.String("from", "", "source currency, required")
    to := flag.String("to", "EUR", "target currency")

    flag.Parse()

    fromCurrency, err := money.ParseCurrency(*from) // (1)
    if err != nil {
        _, _ = fmt.Fprintf(os.Stderr,
            "unable to parse source currency %q: %s.\n",
            *from, err.Error())
        os.Exit(1)
    }

    // TODO：对目标货币重复上述步骤

    // TODO：读取 value 参数，如上文所示

    quantity, err := money.ParseDecimal(value) // (2)
    if err != nil {
        _, _ = fmt.Fprintf(os.Stderr,
            "unable to parse value %q: %s.\n", value, err.Error())
        os.Exit(1)
    }

    amount, err := money.NewAmount( // (3)
        quantity, fromCurrency)
    if err != nil {
        _, _ = fmt.Fprintln(os.Stderr, err.Error())
        os.Exit(1)
    }

    fmt.Println("Amount:", amount, // (4)
        "; Currency:", toCurrency) // (4)
}
\end{gocode}

\begin{codeannotations}
\item 解析货币
\item 解析数量
\item 把值转换为具有相应货币的金额
\item 快速显示结果
\end{codeannotations}

运行代码，享受这场演出。你应该会看到一些类似下面的乱码：

\begin{shellcode}[]
> go run . -from EUR -to CHF 10.50
Amount: {{10 50 2} {EUR 2}}; Currency: {CHF 2}
\end{shellcode}

% SOURCE_PDF_PAGE: 188
对于不了解我们所用结构的人来说，这很难理解。因此，库最好体贴地让自己的几个类型实现 \texttt{Stringer} 接口。

\subsection{Stringer}

查看 \texttt{fmt} 包，可以找到一个很有用的接口：\texttt{Stringer}。\texttt{fmt} 包中的所有格式化函数和打印函数都认识它。它遵循一种 Go 模式：只有一个方法的接口，会采用这个方法名再加上 \texttt{-er} 来命名，例如 \texttt{Reader}、\texttt{Writer} 等等。来看看它的定义：

\begin{gocode}[]
type Stringer interface {
    String() string
}
\end{gocode}

任何拥有 \texttt{String} 方法的类型都实现了 \texttt{fmt.Stringer}，这个方法定义了该值的“原生”格式。只要把值作为操作数传给任何接受字符串的格式，或者传给 \texttt{Print} 这样的无格式打印函数，就会使用 \texttt{String} 方法。在 Go 中，要实现一个接口，类型只需要附带正确的方法即可。

\begin{gocode}[title={代码清单 6.29 \texttt{currency.go}：实现 \texttt{Currency} 的 \texttt{Stringer}}]
// String 实现 Stringer。
func (c Currency) String() string {
    return c.code
}
\end{gocode}

神奇的是，你的 \texttt{Currency} 现在成了一个 \texttt{Stringer}。任何人调用打印函数并把货币作为参数时，都会调用这个方法。试试看：

\begin{shellcode}[]
> go run . -from EUR -to CHF 10.50
Amount: {{10 50 2} EUR}; Currency: CHF
\end{shellcode}

% SOURCE_PDF_PAGE: 189
目标货币现在已经很容易读了。我们来对 \texttt{Decimal} 做同样的事情。\texttt{Decimal} 类型上已经有了一个方法，而 \texttt{simplify} 方法使用指针接收者。Go 不喜欢同一类型的方法既有指针接收者，又有非指针接收者，所以让 \texttt{String()} 接受指针接收者吧。\texttt{simplify} 需要指针接收者。

这里，我们选择使用双重格式化技巧：先用精度创建格式字符串，再用这个字符串格式化数字。例如，精度为两位时，\texttt{format} 变量将是 \texttt{\%d.\%02d}，它会用零填充小数部分。当货币有“分”时，我们不想打印 12.5，而想打印 12.50。这个末尾的 0 正是通过 \texttt{\%02} 格式字符串中的填充添加的。

然而，并非所有货币都有两位精度，我们必须利用货币精度构造这个 \texttt{\%02} 字符串。为此，可以使用负责字符串转换的包所提供的一个函数，这个包的名字很贴切，叫作 \texttt{strconv}。我们使用的函数是 \texttt{strconv.Itoa}，可以把它理解成 \texttt{strconv.Atoi} 的反向操作：

\begin{gocode}[]
decimalFormat := "%d.%0" + strconv.Itoa(int(d.precision)) + "d"
\end{gocode}

我们立刻注意到，精度为 0 是一种边界情况。由于这是一个输出简单的简单测试，我们会在 \texttt{String()} 函数开头检查这个场景。

现在，编写 \texttt{Decimal} 类型的 \texttt{Stringer} 接口实现所需的所有材料都齐了。\texttt{pow10} 的输出给出了一个货币单位中子单位的数量，这意味着只要除以子单位数量，就能取得小数部分和整数部分。最后，可以用格式字符串 \texttt{decimalFormat} 返回打印后的输出。

% SOURCE_PDF_PAGE: 190
\begin{gocode}[title={代码清单 6.30 \texttt{decimal.go}：实现 \texttt{Decimal} 的 \texttt{Stringer}}]
// String 实现 stringer，返回格式化后的 Decimal：
// 若干数字，以及可选的小数点和小数点后的数字。
func (d *Decimal) String() string {
    // 简单情况，无须计算。
    if d.precision == 0 {
        return fmt.Sprintf("%d", d.subunits) // (1)
    }

    centsPerUnit := pow10(d.precision) // (2)
    frac := d.subunits % centsPerUnit
    integer := d.subunits / centsPerUnit

    decimalFormat := "%d.%0" + // (3)
        strconv.Itoa(int(d.precision)) + "d" // (3)
    return fmt.Sprintf(decimalFormat, integer, frac)

}
\end{gocode}

\begin{codeannotations}
\item 没有小数部分时提前退出
\item 把“2”转换为“100”，或把“3”转换为“1000”
\item 创建格式字符串“\texttt{\%d.\%0\{2\}d}”
\end{codeannotations}

即使可以说 \texttt{Currency} 的 \texttt{String} 方法不写单元测试也行，这个方法也需要测试。花一分钟编写测试并检查覆盖率。请记住，覆盖率不会检查你是否得到了充分覆盖，因为即便覆盖率完美，也仍然可能漏掉大量用例；它只是告诉你哪些地方没有覆盖，然后由你判断是否值得花力气扩大覆盖范围。

最后，\texttt{Amount} 也应该实现 \texttt{Stringer} 接口。这可以适配不同语言的标准，但在下面的代码清单中，我们选择 \texttt{22.368 KWD} 作为输出格式。

\begin{gocode}[title={代码清单 6.31 \texttt{amount.go}：实现 \texttt{Amount} 的 \texttt{Stringer}}]
// String 实现 stringer。
func (a Amount) String() string {
    return a.number.String() + " " + a.currency.code
}
\end{gocode}

现在的输出是不是更容易读了？你觉得怎么样：

\begin{shellcode}[]
> go run . -from EUR -to CHF 10.50
10.50 EUR CHF
\end{shellcode}

% SOURCE_PDF_PAGE: 191
既然已经实现了所有 \texttt{Stringer} 代码，现在可以调用 \texttt{Convert} 函数了。

\subsection{Convert}

\texttt{main} 函数唯一剩下的事情就是调用 \texttt{Convert}。下面代码清单中的代码可以编译运行——它真漂亮。

\begin{gocode}[title={代码清单 6.32 \texttt{main.go}：\texttt{main} 函数的结尾}]
func main() {
    // ...

    convertedAmount, err := money.Convert(amount, toCurrency)
    if err != nil {
        _, _ = fmt.Fprintf(os.Stderr,
            "unable to convert %s to %s: %s.\n",
            amount, toCurrency, err.Error())
        os.Exit(1)
    }

    fmt.Printf("%s = %s\n", amount, convertedAmount))
}
\end{gocode}

不过，这里还有最后一个问题必须解决。尽管进行了大量测试，我们还是漏掉了一个相当明显的问题——你发现这个 bug 了吗？如果试着运行过工具，可能已经注意到了。当传入金额的小数位数少于货币精度时，显示该金额用的是输入数字的位数，而不是其货币的位数！下面是一个例子：

\begin{shellcode}[]
> go run . -from BHD -to CHF 12.5
12.5 BHD = 25.00 CHF
\end{shellcode}

查看 \texttt{ParseCurrency} 代码中的 \texttt{switch}，可以看到 1 巴林第纳尔等于 1,000 费尔（fulūs），因此应该写成 \texttt{12.500 BHD = 25.00 CHF}，让第纳尔的小数点后有三位数字。

问题的根源在 \texttt{NewAmount} 函数中。我们来把货币精度考虑进去，从而修复它
